\chapter{Análisis de Impacto}

Con el desarrollo del proyecto se han identificado distintos contextos donde tiene un impacto. Por un lado, desde el punto de vista personal, el trabajo permitió adquirir experiencia práctica en el análisis, diseño, desarrollo, despliegue y mantenimiento de una aplicación software completa. Además, se tomaron decisiones técnicas para conseguir los objetivos de la forma más óptima posible. Por otro lado, desde el punto de vista empresarial, el proyecto demuestra la viabilidad de construir una plataforma de streaming de contenido en línea. Se han analizado proveedores, tecnologías y herramientas teniendo en cuenta los costos y además, se crearon campañas publicitarias para atraer a nuevos usuarios. 

Desde un punto de vista social, una plataforma de este tipo facilita la difusión de contenido y la creación de comunidades. No obstante, también puede generar efectos sociales adversos como la distribución de contenido dañino o inapropiado, lo que supuso crear métodos de moderación.  

En cuanto al contexto económico, el uso de servicios externos en la nube y proveedores permite reducir la inversión inicial. Sin embargo, esto introduce costos variables que aumentan con el volumen de usuarios en la plataforma. Por esta razón, los servicios y proveedores fueron analizados tomando esto en cuenta junto con su complejidad y escalabilidad. 

El proyecto tiene un impacto medioambiental debido al consumo energético asociado al procesamiento, almacenamiento, distribución de videos y a las operaciones de la plataforma. La incorporación de límites de uso para los usuarios reduce la utilización de recursos de la aplicación y por tanto, su consumo energético. 

Finalmente, según los Objetivos de Desarrollo Sostenible, el proyecto se alinea con ODS 4 (Educación de calidad) por su potencial para difundir contenido educativo; ODS 8 (Trabajo decente y crecimiento) ya que se crean oportunidades laborales para los creadores de contenido quienes son remunerados por cargar contenido a la plataforma; ODS 9 (Industria, innovación e infraestructura) al desarrollar una solución tecnológica con herramientas y servicios modernos. 



    % En este capítulo se realizará un análisis del impacto potencial de los resultados obtenidos durante la realización del trabajo, en los diferentes contextos para los que se aplique:\footnote{Téngase en cuenta que no se espera que un trabajo tenga impacto en todos los contextos}
                        
    % \begin{itemize}
    % \item Personal
    % \item Empresarial
    % \item Social
    % \item Económico
    % \item Medioambiental
    % \item Cultural
    % \end{itemize}

    % En dicho análisis se destacarán los beneficios esperados, así como también los posibles \textbf{efectos adversos}.
    % Además, se harán notar aquellas decisiones tomadas a lo largo del trabajo que tienen como base la consideración del impacto.

    % Se recomienda analizar también el potencial impacto respecto a los Objetivos de Desarrollo Sostenible (ODS), de la Agenda 2030, que sean relevantes para el trabajo realizado (\href{https://www.un.org/sustainabledevelopment/es/objetivos-de-desarrollo-sostenible/}{ver enlace})

